\documentclass[11pt]{article}
\usepackage[margin=1in]{geometry}
\usepackage{setspace}
\usepackage{titlesec}
\usepackage{parskip}
\usepackage{hyperref}

\titleformat{\section}{\large\bfseries}{\thesection}{1em}{}
\renewcommand{\thesection}{\arabic{section}.}

\title{\vspace{-2cm}Proposal for Pursuing a PhD in Systems Engineering}
\author{}
\date{\today}

\begin{document}

\begin{flushright}
\textbf{To:} Krystal Schwartz, IT Management \\
\textbf{CC:} Rajesh Yadav, IT Management \\
\textbf{From:} Stella Dunn, IT Engineer, Intermediate \\
\textbf{Date:} \today
\end{flushright}

\section{Background}

Since joining the company on July 1, 2024, I have led a complete overhaul of the Vulnerability Management team's reporting processes. At the time, the system consisted of a disjointed collection of scripts that frequently failed and required approximately five hours of manual work from two team members each week. Over the past 10 months, I have transformed this into a streamlined, reliable system with dramatic results:

\begin{itemize}
    \item \textbf{Unified the codebase} into a maintainable structure with a shared library and installation process.
    \item \textbf{Built a command-line interface (CLI)} to simplify user interaction, with plans for a full graphical interface underway.
    \item \textbf{Implemented robust error handling} and added new functionality that eliminated multiple manual workflows.
    \item \textbf{Reduced total processing time} from roughly \textbf{5 hours to just 30 minutes end-to-end}, with minimal manual intervention.
    \item \textbf{Integrated external APIs} to enhance data quality and broaden the system's capabilities.
    \item \textbf{Developed supporting tools}, including a mailer script and an ad hoc enrichment utility, to support evolving team needs.
\end{itemize}

In parallel, I co-led a key systems recovery project that involved automating the programming of another automation platform to bring thousands of disconnected agents back online. This effort required not only technical precision, but also strategic coordination across tools and teams.

Additionally, I temporarily led the team of U.S.-based ITAR administrators for the PSF tool, ensuring continuity and operational efficiency while permanent personnel were identified and onboarded. This role required cross-team coordination and proactive issue resolution to maintain compliance and access integrity during a transitional period.

These initiatives reflect a high degree of systems-level thinking, which is foundational to both my current role and the field of Systems Engineering. Pursuing a PhD in this discipline will expand my ability to design, manage, and optimize complex technical ecosystems—directly benefiting our team and broader organizational goals.

\section{Program Overview}

Colorado State University (CSU) offers a fully online PhD in Systems Engineering tailored for experienced professionals. The program requires 72 credits (with a 42-credit option for those with a relevant master's degree), is asynchronous-friendly, and culminates in original research and a dissertation.

\begin{itemize}
    \item \textbf{Format:} 100\% online (synchronous and asynchronous components).
    \item \textbf{Start Date:} Spring semester (application deadline: November 1).
    \item \textbf{Credit Requirement:} 72 credits at \$1,175 per credit; possible reduction to 42 credits (~\$49,350) based on transfer credit approval.
    \item \textbf{Content Areas:} Systems architecture, risk analysis, simulation, machine learning, data mining, cybersecurity, and software-intensive systems.
\end{itemize}

\section{Alignment with Company Goals}

My responsibilities already include engineering automation, managing complex systems, and integrating data from multiple sources. The PhD program directly complements these tasks by deepening my knowledge of:

\begin{itemize}
    \item Advanced systems thinking and architecture.
    \item Artificial intelligence and machine learning.
    \item Cybersecurity and information assurance.
    \item Large-scale data mining and analytics.
\end{itemize}

CSU's program emphasizes solving real-world, multidisciplinary problems and includes research areas that align with our department's transition to a software-engineering focus. This education will increase my capacity to innovate, lead complex technical initiatives, and contribute to the company's strategic goals.

\section{Benefits to the Organization}

\begin{itemize}
    \item \textbf{Cybersecurity Innovation:} I will bring back applied research knowledge in information assurance and threat mitigation.
    \item \textbf{Process Automation:} Enhanced skills in AI and data systems will allow for more advanced, efficient reporting and automation tools.
    \item \textbf{Leadership Development:} I will grow into a stronger technical leader capable of mentoring others and driving strategic initiatives.
    \item \textbf{Retention and Morale:} Supporting this degree signals a strong investment in employee development.
\end{itemize}

\section{Time Management Plan}

The PhD program is structured for working professionals and allows coursework to be scheduled outside of core working hours. I plan to:

\begin{itemize}
    \item Complete classes and study sessions during evenings and weekends.
    \item Maintain full-time work status and continue meeting all performance expectations.
    \item Utilize my current flexible schedule only as needed, and communicate any short-term adjustments well in advance.
    \item Provide regular status updates on both academic and work progress.
\end{itemize}

\section{Financial Consideration}

\begin{itemize}
    \item \textbf{Tuition:} \$1,175 per credit hour; estimated total cost is between \$49,350 and \$84,600 depending on approved transfer credits.
    \item \textbf{Additional Costs:} University fees and textbooks (estimates pending).
    \item \textbf{Funding Request:} I request full coverage of tuition and fees under the company's Educational Assistance Program, with the understanding that final cost will be reduced by any transfer credits, scholarships, or grants I may receive.
\end{itemize}

As outlined in the Educational Assistance Program guidelines, financial support is discretionary and contingent on program relevance and funding availability. My proposed field of study, Systems Engineering, falls under Science and Technology Studies and qualifies for support under the current policy framework.

\section{Next Steps}

\begin{enumerate}
    \item Request transcript evaluation from CSU to confirm credit transfer eligibility.
    \item Identify and confirm a faculty advisor for the PhD application.
    \item Submit application by November 1 for Spring semester admission.
    \item Apply for scholarships and university aid.
    \item Coordinate with HR to submit formal tuition assistance documentation.
    \item Begin coursework upon admission and provide periodic progress reports.
\end{enumerate}

\section*{Conclusion}

I am committed to continuing to deliver high-value work to our team while advancing my education. This PhD in Systems Engineering is directly aligned with my role, department needs, and the company's strategic direction. I respectfully request your approval to proceed with the application and funding request process.

\vspace{1cm}
\noindent Sincerely, \\
Stella Dunn \\
IT Engineer, Intermediate

\end{document}
